% Generated draft: requires mathematical and source review.
\documentclass[11pt]{article}
\usepackage[margin=25mm]{geometry}
\usepackage{fontspec}
\setmainfont{Latin Modern Roman}
\usepackage{amsmath,amssymb,mathtools,unicode-math}
\setmathfont{Latin Modern Math}
\usepackage{enumitem,needspace,xurl,hyperref}
\setlength{\parindent}{0pt}
\setlength{\parskip}{5pt}
\setlength{\emergencystretch}{3em}
\newcommand{\N}{\mathbb{N}}
\newcommand{\Z}{\mathbb{Z}}
\newcommand{\Q}{\mathbb{Q}}
\newcommand{\R}{\mathbb{R}}
\newcommand{\C}{\mathbb{C}}
\newcommand{\F}{\mathbb{F}}
\newcommand{\A}{\mathbb{A}}
\newcommand{\PP}{\mathbb{P}}
\newcommand{\E}{\mathbb{E}}
\newcommand{\eps}{\varepsilon}
\newcommand{\dd}{\,\mathrm{d}}
\begin{document}
\section*{UnsolvedMath ID: 800301}
\textbf{Source problem code:} catalogue entry
\subsection{Definitions and mathematical statement}
Let \(S^1=\R/\Z\), and let \(m\) be normalized Lebesgue measure on \(S^1\). Write \(\operatorname{Diff}^{2}_{+}(S^1)\) for the group of orientation-preserving \(C^2\) circle diffeomorphisms. A probability distribution on a finitely generated subgroup \(G\leq\operatorname{Diff}^{2}_{+}(S^1)\) is a function \(\mu:G\to[0,1]\) satisfying \(\sum_{g\in G}\mu(g)=1\). A Borel probability measure \(\nu\) on \(S^1\) is \(\mu\)-stationary if
\[
\nu=\sum_{g\in G}\mu(g)\,g_*\nu,
\qquad (g_*\nu)(A)=\nu(g^{-1}(A))
\]
for every Borel set \(A\subseteq S^1\). Absolute continuity \(\nu\ll m\) means that \(m(A)=0\) implies \(\nu(A)=0\) for every Borel set \(A\).

The supplied question asks whether, for every finitely generated \(G\leq\operatorname{Diff}^{2}_{+}(S^1)\), there exist a non-degenerate probability distribution \(\mu\) on \(G\) and a Borel probability measure \(\nu\) on \(S^1\) such that
\[
\nu=\sum_{g\in G}\mu(g)\,g_*\nu
\quad\text{and}\quad
\nu\ll m.
\]
The source record does not define ``non-degenerate.'' In particular, it does not specify whether the support must generate \(G\) as a group, generate it as a semigroup, or equal \(G\). This qualification remains unresolved. The displayed formulation interprets ``the stationary measure'' existentially; the record supplies neither a uniqueness assertion nor hypotheses guaranteeing uniqueness. No further restrictions on the action are supplied.
\subsection{Short English statement}
Navas asks whether every finitely generated group of orientation-preserving \(C^2\) circle diffeomorphisms admits a non-degenerate random walk having an absolutely continuous stationary probability measure. The meaning of ``non-degenerate'' and the intended uniqueness qualification require clarification from the source.
\subsection{Sources}
\begin{itemize}
\item Supplied local UnsolvedMath registry record: permanent ID \(800301\), exact source problem code \texttt{catalogue entry}, title ``Group Actions on 1-Manifolds --- Question 18.'' Attribution: Navas, proposed in 2017. The registry labels the problem open; this is a supplied status, not an independently verified mathematical or literature assessment.
\item Navas, \emph{Group Actions on 1-Manifolds: A List of Very Concrete Open Questions} (2017), \url{https://arxiv.org/abs/1712.06462}. The supplied provenance identifies the item as ``Question 18 (\#17),'' records access on 2026-10-16, and gives extraction method \texttt{arxiv\_latex}. These are supplied provenance details; the reference has not been independently consulted for this draft.
\item Supplied upstream revision: \texttt{d371ee917ac9647fbb06cd6fc529deec6f3dda82}.
\end{itemize}
\end{document}
